% Informe de evidencias de ejecución — Avance 1 (Consultor Normativo)
% Estilo inspirado en el informe LaTeX del prototipo BetaPRelectiva1.
%
% CÓMO COMPILAR
% --------------
% Opción A (sin instalar nada): subir esta carpeta a https://overleaf.com
%   (New Project > Upload Project, o pegar este archivo en un proyecto nuevo).
% Opción B (local, con MiKTeX o TeX Live instalado):
%   pdflatex informe_avance1.tex
%   pdflatex informe_avance1.tex   (segunda pasada para el índice/numeración)
%
% CAPTURAS DE PANTALLA
% ---------------------
% Coloque las imágenes en ../img/ con estos nombres exactos:
%   evidencia_00.png ... evidencia_07.png  (y opcionalmente evidencia_doc.png)
% Si el archivo existe, se incluye automáticamente; si no, se muestra un
% recuadro de marcador de posición para que sepa qué falta.

\documentclass[11pt]{article}

\usepackage[spanish]{babel}
\usepackage[utf8]{inputenc}
\usepackage[T1]{fontenc}
\usepackage[margin=2.5cm]{geometry}
\usepackage{graphicx}
\usepackage{hyperref}
\usepackage{enumitem}
\usepackage{xcolor}
\usepackage{listings}

\lstset{
  basicstyle=\ttfamily\small,
  breaklines=true,
  columns=fullflexible,
  backgroundcolor=\color{black!5},
  frame=single,
  framesep=4pt,
}

\newcommand{\evidenciaimg}[3]{%
  \begin{figure}[h!]
    \centering
    \IfFileExists{#1}{%
      \includegraphics[width=0.85\textwidth]{#1}%
    }{%
      \fbox{\begin{minipage}[c][5.5cm][c]{0.85\textwidth}
        \centering
        \textbf{[Espacio para captura de pantalla]}\\[6pt]
        \texttt{\detokenize{#1}}
      \end{minipage}}%
    }
    \caption{#2}
    \label{#3}
  \end{figure}
}

\title{\textbf{Evidencia de Ejecución y Análisis del Prompt}\\
  \large Consultor Normativo — Asistente de Orientación sobre Normativa Colombiana}
\author{Joan Sebastian Lara Fuenmayor}
\date{\today}

\begin{document}

\maketitle

\section{Diseño y Estructuración del Prompt}

El asistente «Consultor Normativo» se construyó siguiendo tres técnicas de
ingeniería de prompts vistas en clase: configuración del \emph{system prompt},
\emph{few-shot prompting} y delimitadores. El objetivo es que el modelo
responda dudas normativas de una persona sin formación jurídica, citando
siempre la norma y el artículo exactos, y reconociendo explícitamente cuándo
la información \textbf{no} está en sus fuentes (control de alucinaciones).

\subsection{System Prompt}

Se define a la IA con el rol de \textbf{«Consultor Normativo»}, no abogado,
cuyo conocimiento se limita EXCLUSIVAMENTE al texto entregado dentro de la
etiqueta \texttt{<contexto\_normativo>} de cada mensaje. El system prompt
incluye:
\begin{itemize}
  \item \textbf{Dominio acotado:} Ley 1581 de 2012, Ley 820 de 2003,
    Ley 1480 de 2011, Código Civil, CPACA (Ley 1437 de 2011), Estatuto
    Tributario y Código General del Proceso (Ley 1564 de 2012).
  \item \textbf{Reglas estrictas (zero-hallucination):} prohibido inventar
    números de artículo, nombres de normas o citas textuales; si la consulta
    es jurídica pero el contexto no la cubre, la respuesta debe ser
    exactamente: \emph{"La información solicitada no se encuentra en las
    fuentes consultadas."}
  \item \textbf{Lógica condicional} (equivalente a \texttt{if/else} en
    lenguaje natural) para distinguir entre: consulta resuelta por el
    contexto, consulta jurídica sin fundamento en el contexto, y consulta
    fuera del dominio del asistente.
  \item \textbf{Formato de salida forzado:} JSON estructurado (o Markdown,
    según configuración), sin texto adicional.
\end{itemize}

\subsection{Few-Shot Prompting}

Se incluyen 5 ejemplos ya resueltos, con el mismo esquema JSON, cubriendo
casos variados: arrendamiento, datos personales, consumidor, una consulta
\textbf{fuera de alcance} y una consulta \textbf{sin fundamento} en el
contexto. Esto fija el formato de salida y el criterio de respuesta antes de
recibir la consulta real del usuario.

\subsection{Delimitadores}

Se usan etiquetas tipo XML para separar de forma inequívoca cada bloque del
prompt de usuario:

\begin{lstlisting}
<contexto_normativo> ...textos de las leyes... </contexto_normativo>
<consulta_usuario>   ...la pregunta del usuario... </consulta_usuario>
<instrucciones>      ...qué debe hacer el modelo... </instrucciones>
\end{lstlisting}

Se prefirieron sobre la triple comilla o la triple tilde invertida
(\texttt{```}) porque permiten anidar bloques (cada ejemplo few-shot lleva
sus propias secciones internas) y reducen el riesgo de que la consulta del
usuario se haga pasar por una instrucción (inyección de prompt).

\subsection{Formato de salida}

\begin{lstlisting}
{
  "respuesta": "explicacion clara para una persona sin formacion juridica",
  "fundamento_normativo": [
    { "norma": "Ley 1581 de 2012", "articulo": "Articulo 8",
      "cita_textual": "..." }
  ],
  "nivel_confianza": "alto | medio | bajo",
  "advertencia": "Este contenido es orientacion informativa...",
  "requiere_abogado": true
}
\end{lstlisting}

\section{Entorno usado}

\begin{center}
\begin{tabular}{|l|l|}
\hline
\textbf{Dato} & \textbf{Valor} \\
\hline
Motor LLM (\texttt{LLM\_PROVIDER}) & \\
\hline
Modelo & \\
\hline
Sistema operativo & \\
\hline
Fecha de ejecución & \\
\hline
\end{tabular}
\end{center}

\section{Evidencias de Ejecución}

A continuación se presentan las pruebas realizadas en la consola que
demuestran el comportamiento del sistema frente a distintos escenarios.

\subsection{Evidencia 0 — Estructura del prompt (sin llamar al modelo)}

\begin{lstlisting}
python src/main.py --ver-prompt --normas ley_820_2003 --consulta "¿Puedo terminar el contrato de arriendo antes de tiempo?"
\end{lstlisting}

\textbf{Qué observar:} el System Prompt con las secciones \texttt{<dominio>},
\texttt{<reglas>}, \texttt{<estilo>} y \texttt{<formato\_salida>}; y el User
Prompt con \texttt{<ejemplos>} (few-shot), \texttt{<contexto\_normativo>},
\texttt{<consulta\_usuario>} e \texttt{<instrucciones>}.

\evidenciaimg{../img/evidencia_00.png}{Estructura ensamblada del prompt.}{fig:ev0}

\subsection{Evidencia 1 — Consulta CON fundamento (arrendamiento)}

\begin{lstlisting}
python src/main.py --normas ley_820_2003
# Consulta> Si quiero terminar mi contrato de arriendo antes de tiempo, ¿con cuanta anticipacion debo avisar y tengo que pagar algo?
\end{lstlisting}

\textbf{Qué observar:} salida en JSON válido; \texttt{fundamento\_normativo}
cita la \textbf{Ley 820 de 2003, Artículo 24} con \texttt{cita\_textual}
tomada del contexto; \texttt{nivel\_confianza: "alto"}.

\evidenciaimg{../img/evidencia_01.png}{Consulta con fundamento normativo — arrendamiento.}{fig:ev1}

\subsection{Evidencia 2 — Consulta CON fundamento (datos personales)}

\begin{lstlisting}
python src/main.py --normas ley_1581_2012
# Consulta> ¿Que derechos tengo sobre mis datos personales y en cuanto tiempo deben responderme una consulta?
\end{lstlisting}

\textbf{Qué observar:} cita \textbf{Ley 1581 de 2012, Artículo 8} y
\textbf{Artículo 14} (término de 10 días hábiles). El modelo no agrega
información que no esté en el contexto.

\evidenciaimg{../img/evidencia_02.png}{Consulta con fundamento normativo — datos personales.}{fig:ev2}

\subsection{Evidencia 3 — Consulta CON fundamento (consumidor / retracto)}

\begin{lstlisting}
python src/main.py --normas ley_1480_2011
# Consulta> Compre un electrodomestico por internet y me arrepenti al dia siguiente. ¿Puedo devolverlo?
\end{lstlisting}

\textbf{Qué observar:} cita \textbf{Ley 1480 de 2011, Artículo 47}; menciona
los \textbf{5 días hábiles} y la devolución del dinero.

\evidenciaimg{../img/evidencia_03.png}{Consulta con fundamento normativo — derecho de retracto.}{fig:ev3}

\subsection{Evidencia 4 — Consulta FUERA DE ALCANCE (guardrail)}

\begin{lstlisting}
python src/main.py
# Consulta> ¿Me recomiendas invertir en criptomonedas este mes?
\end{lstlisting}

\textbf{Qué observar:} el asistente \textbf{no responde} la pregunta
financiera; redirige e indica su alcance. \texttt{fundamento\_normativo: []},
\texttt{nivel\_confianza: "bajo"}.

\evidenciaimg{../img/evidencia_04.png}{Guardrail ante consulta fuera del dominio.}{fig:ev4}

\subsection{Evidencia 5 — Consulta jurídica SIN fundamento en el contexto}

\begin{lstlisting}
python src/main.py --normas codigo_civil
# Consulta> ¿La ley colombiana permite el matrimonio entre primos hermanos?
\end{lstlisting}

\textbf{Qué observar:} el fragmento de Código Civil cargado trata de
contratos y responsabilidad, no de matrimonio. La respuesta debe ser
exactamente: \emph{"La información solicitada no se encuentra en las fuentes
consultadas."}, con \texttt{fundamento\_normativo: []}. \textbf{Esto
demuestra el control de alucinaciones vía prompt.}

\evidenciaimg{../img/evidencia_05.png}{Control de alucinaciones ante ausencia de fundamento.}{fig:ev5}

\subsection{Evidencia 6 — Mismo prompt, otro motor (opcional)}

Se repite una de las consultas cambiando \texttt{LLM\_PROVIDER} en
\texttt{.env} (de \texttt{ollama} a \texttt{gemini} o viceversa) para
evidenciar que la estructura de prompts es independiente del motor.

\evidenciaimg{../img/evidencia_06.png}{Misma consulta ejecutada en un motor distinto.}{fig:ev6}

\subsection{Evidencia 7 — Formato Markdown (opcional)}

\begin{lstlisting}
python src/main.py --formato markdown --normas estatuto_tributario
# Consulta> ¿Cual es la sancion por presentar tarde una declaracion tributaria?
\end{lstlisting}

\textbf{Qué observar:} el mismo contenido pero renderizado en Markdown, según
lo definido en el bloque \texttt{<formato\_salida>} del System Prompt.

\evidenciaimg{../img/evidencia_07.png}{Salida en formato Markdown.}{fig:ev7}

\section{Conclusión}

% TODO: redactar 3-5 líneas sobre qué se logró con la estructuración del
% prompt, cómo el System Prompt y el Few-Shot fuerzan el formato JSON, y cómo
% los delimitadores y la lógica condicional evitan que el modelo responda
% fuera de las fuentes.

La combinación de un System Prompt fuertemente parametrizado, ejemplos
few-shot consistentes y delimitadores tipo XML demuestra ser eficaz para
tareas de orientación normativa. La delimitación explícita del contexto y la
restricción del formato de salida garantizan la confiabilidad del sistema en
un dominio donde la precisión de la información es crítica, evitando que el
modelo complete con conocimiento propio lo que no está en las fuentes
entregadas.

\appendix
\section{Anexo (opcional): modo de análisis documental}

Como capacidad adicional, el proyecto incorpora un segundo modo,
independiente del Consultor Normativo, para analizar cualquier documento
\texttt{.txt}/\texttt{.pdf} cargado por el usuario (contratos, manuales,
políticas, etc.) con el mismo criterio de cero-alucinación, citando siempre
el fragmento textual exacto que respalda la respuesta:

\begin{lstlisting}
python src/main.py --modo documento --documento contrato.pdf
\end{lstlisting}

\evidenciaimg{../img/evidencia_doc.png}{Modo de análisis documental sobre un documento arbitrario.}{fig:evdoc}

\end{document}
